\documentclass[11pt,a4paper]{article}
\usepackage[margin=1in]{geometry}
\usepackage{amsmath}
\usepackage{amssymb}
\usepackage{algorithm}
\usepackage{algpseudocode}
\usepackage{xcolor}
\usepackage{hyperref}
\usepackage{booktabs}
\usepackage{graphicx}

\title{Automated Conversion of Topologically Optimised Finite Element Mesh-Based Geometries into Feature-Based Parametric CAD Models}
\author{}
\date{}

\newcommand{\fig}[1]{\textit{[Figure: #1]}}
\newcommand{\tbl}[1]{\textit{[Table: #1]}}
\newcommand{\alg}[1]{\textit{[Algorithm: #1]}}

\begin{document}

\maketitle

\section{Introduction}
\label{sec:intro}

\subsection{Motivation}
\label{subsec:motivation}

Topology optimisation (SIMP, level-set, etc.) produces structurally efficient but voxelised density fields that are non-parametric and non-manufacturable. The ``interpretation gap''---the gap between engineers' access to a density cloud and their ability to directly edit, dimension, or manufacture such a cloud---is the central problem motivating this work. Engineers cannot directly edit a density cloud; they need feature-based parametric CAD models. Industry adoption of topology optimisation is limited precisely because this manual interpretation step is time-consuming, subjective, and error-prone. Three-dimensional domains with tens of thousands of elements produce complex, organic geometries that would take weeks to reconstruct manually.

\subsection{Problem Statement}
\label{subsec:problem}

Existing automated approaches produce either beam-only reconstructions (Yin 2020, Polak 2023) or surface-only models (Amroune 2021, Feng 2025)---none handle structures with both beam-like members and plate-like regions simultaneously. Frame optimisation in prior work runs as a single pass with no analysis of convergence behaviour. No complete open-source Python pipeline exists from topology optimisation solver through to parametric CAD output.

The research question is explicit: \emph{Can a fully automated pipeline classify, reconstruct, and optimise mixed beam-plate structures from 3D topology optimisation results into editable parametric CAD models?}

\subsection{Key Contributions}
\label{subsec:contributions}

This dissertation presents four contributions:

\begin{enumerate}
    \item \textbf{Python reimplementation of Yin's pipeline.} A complete Python port of Yin's (2020) MATLAB topology optimisation-to-CAD pipeline with boundary condition tag propagation throughout all pipeline stages.

    \item \textbf{Multi-signal zone classification.} A hybrid zone classification method combining PCA shape analysis, aspect ratio scoring, boundary condition tag density, and Euclidean Distance Transform uniformity to automatically identify beam-like vs.\ plate-like regions.

    \item \textbf{Iterative optimisation convergence.} Demonstration of 2-iteration convergence in iterative size-then-layout frame optimisation---an empirical finding not reported in prior work, with practical implications for computational cost.

    \item \textbf{End-to-end automation.} Complete automation from Top3D SIMP topology optimisation to parametric FreeCAD models, including cubic Bézier curved beam fitting and crash-resistant macro generation.
\end{enumerate}

\fig{1.1: Overview diagram---topology optimisation density field (left) $\rightarrow$ ``This Work'' arrow $\rightarrow$ parametric CAD model (right). Use actual $150 \times 40 \times 4$ cantilever pipeline output.}

\section{Review of Related Literature}
\label{sec:litreview}

\subsection{Topology Optimisation for Structural Design}
\label{subsec:simp}

The Solid Isotropic Material with Penalisation (SIMP) method (Bendsoe \& Kikuchi 1988; Sigmund 2001) is the foundation of modern density-based topology optimisation. The material interpolation law is
\begin{equation}
    E(\rho) = E_{\min} + \rho^p (E_0 - E_{\min}),
\end{equation}
where $\rho$ is material density (0 to 1), $p$ is the penalisation power (typically 3), and $E_0$ and $E_{\min}$ are the solid and void elastic moduli.

The 88-line 3D extension (Top3D) is the starting point for this work. Sensitivity analysis computes material derivatives, density filtering prevents checkerboarding, and the Optimality Criteria (OC) update scheme with bisection satisfies the volume constraint. Topology optimisation produces optimal material distributions, but the output is a density field---a continuous scalar field over the domain, not a geometry. The interpretation of this density field into manufacturable shapes is the subject of this dissertation.

\subsection{Skeletonisation and Medial Axis Approaches}
\label{subsec:skeletonisation}

Digital topology foundations (Bertrand \& Malandain) establish the theory of 26-connectivity (foreground) and 6-connectivity (background) on the integer voxel lattice. A \emph{simple point} is a voxel whose removal does not change the topology of either foreground or background. The medial axis (or skeleton) is the locus of centres of maximal inscribed balls.

Parallel directional thinning is an algorithm that iteratively removes border voxels in six directions (+X, $-$X, +Y, $-$Y, +Z, $-$Z) while preserving topology via simple point checking and end voxel preservation. Yin's Algorithm 3.1 is a canonical implementation. The Euclidean Distance Transform (EDT) approximates medial axis thickness at each voxel: $\text{radius} = 2 \times \text{EDT}_{\text{value}} \times \text{pitch}$.

\subsection{Frame-Based CAD Reconstruction}
\label{subsec:frame_methods}

\subsubsection{Yin (2020)}

Yin's approach: SIMP topology optimisation $\rightarrow$ thresholding $\rightarrow$ homotopic skeletonisation (Algorithm 3.1) $\rightarrow$ graph extraction (Algorithm 4.1) $\rightarrow$ size and layout optimisation $\rightarrow$ Constructive Solid Geometry (CSG) CAD model (cylinders + spheres).

\textbf{Strengths:} Topologically sound thinning via simple point checking, boundary condition (BC) protection during thinning, complete frame extraction and optimisation pipeline.

\textbf{Limitations:} Beam-only reconstruction (no plates), single optimisation pass (no convergence analysis), MATLAB implementation (not open-source), no curved beam representation.

\subsubsection{Polak (2023)}

Biomimetic topology optimisation $\rightarrow$ skeletonisation $\rightarrow$ evolutionary refinement loop $\rightarrow$ FreeCAD export with cylinders and spheres.

\textbf{Strengths:} Iterative refinement concept, FreeCAD automation.

\textbf{Limitations:} Beam-only reconstruction, heuristic-based evolutionary refinement (no formal optimisation formulation), no convergence analysis, limited to simple geometries.

\subsection{Surface and Lofting-Based Methods}
\label{subsec:surface_methods}

\subsubsection{Amroune (2021)}

Curve-skeletonisation $\rightarrow$ cross-section extraction at regular intervals $\rightarrow$ lofting surfaces $\rightarrow$ junction blending $\rightarrow$ solid CAD model with finite element analysis (FEA) validation.

\textbf{Strengths:} Handles varying cross-sections along branches, FEA validation included, smooth surface output.

\textbf{Limitations:} Requires clear beam-like topology, no plate handling, no structural re-optimisation after reconstruction.

\subsubsection{Feng (2025)}

Mesh smoothing $\rightarrow$ skeleton extraction $\rightarrow$ Rotation Minimising Frames (RMF) for consistent orientation $\rightarrow$ lofting surfaces $\rightarrow$ Boolean union $\rightarrow$ editable CAD model.

\textbf{Strengths:} Smooth continuous surfaces, produces editable STEP files for commercial CAD.

\textbf{Limitations:} RMF computation is complex, no structural re-optimisation after reconstruction, no mixed element types.

\subsection{Summary and Research Gap}
\label{subsec:gap}

No prior work handles hybrid beam-plate structures within a single automated pipeline. No prior work demonstrates or analyses iterative frame optimisation convergence behaviour. No fully open-source Python implementation exists for this class of problem.

This work extends Yin (2020) with:
\begin{itemize}
    \item Hybrid beam-plate zone classification via multi-signal scoring,
    \item Iterative size-then-layout optimisation with demonstrated convergence analysis,
    \item Complete Python implementation with open-source distribution,
    \item FreeCAD automation and curved beam representation via cubic Bézier curves.
\end{itemize}

\tbl{1.1: Comparison of 4 reference papers across: Language, Input format, Element types handled, Optimisation stages, CAD output format, Plate handling (Y/N), Iterative optimisation (Y/N), Open source (Y/N)}

\section{Methodology}
\label{sec:methodology}

\subsection{Top3D Topology Optimisation}
\label{subsec:top3d}

\subsubsection{SIMP Formulation}

Material interpolation follows equation \eqref{eq:simp} with penalisation power $p = 3$. Each element uses an 8-node hexahedral (H8) brick element with 24 degrees of freedom (DOF per element). The element stiffness matrix $\mathbf{K}_E$ is computed via analytical integration.

The optimisation problem is:
\begin{align}
    \min_{\boldsymbol{\rho}} \quad & C = \mathbf{F}^T \mathbf{u} \\
    \text{subject to} \quad & \sum_{e} \rho_e v_e \leq V^*, \quad 0 < \rho_{\min} \leq \rho_e \leq 1,
\end{align}
where $C$ is compliance, $\mathbf{F}$ is the load vector, $\mathbf{u}$ is the displacement vector (from $\mathbf{K} \mathbf{u} = \mathbf{F}$), $\rho_e$ is element density, $v_e$ is element volume, and $V^*$ is the target volume.

\subsubsection{Sensitivity Analysis and Filtering}

The sensitivity of compliance with respect to element density is
\begin{equation}
    \frac{\partial C}{\partial \rho_e} = -p (E_0 - E_{\min}) \rho_e^{p-1} \mathbf{u}_e^T \mathbf{K}_E \mathbf{u}_e.
\end{equation}

A density filter with radius $r_{\min}$ applies a weighted average of neighbourhood sensitivities to prevent checkerboarding:
\begin{equation}
    \bar{\psi}_e = \frac{\sum_f H_{ef} \psi_f}{\sum_f H_{ef}}, \quad H_{ef} = \max(0, r_{\min} - \text{dist}(e,f)).
\end{equation}

\subsubsection{Optimality Criteria Update and Convergence}

The OC update rule with move limit (typically 0.2) is applied. The Lagrange multiplier for the volume constraint is found via bisection to satisfy the constraint. The convergence criterion is
\begin{equation}
    \max_e |\rho_e^{\text{new}} - \rho_e^{\text{old}}| < 0.01.
\end{equation}

\subsubsection{Boundary Condition Tag Generation}

This is a critical addition not in the original 88-line code. Elements adjacent to fixed-displacement nodes are tagged as type 1 (support); elements adjacent to loaded nodes are tagged as type 2 (load). These tags propagate through thinning, graph extraction, and into optimisation, ensuring BC points survive all pipeline stages. Without tags, thinning can remove load introduction points, making the reconstructed frame structurally meaningless.

\tbl{2.1: Top3D default parameters (nelx, nely, nelz, volfrac, penal, rmin, max\_loop, vol\_thresh)}

\subsection{Skeletonisation}
\label{subsec:skel}

\subsubsection{Thresholding and Binarisation}

The continuous density field is converted to a binary voxel grid via
\begin{equation}
    \text{solid} = (\rho > \text{vol\_thresh}),
\end{equation}
with default threshold 0.3 (below the typical SIMP penalisation cutoff). Threshold sensitivity is important: too high loses structural members, too low includes noise.

\subsubsection{Parallel Directional Thinning (Algorithm 1)}

Yin's Algorithm 3.1 adapted with BC protection:

\alg{1: Pseudocode for parallel directional thinning with BC protection}

For each iteration:
\begin{enumerate}
    \item \textbf{For each direction} in $\{+X, -X, +Y, -Y, +Z, -Z\}$:
    \begin{itemize}
        \item Identify border voxels: solid voxels adjacent to background in that direction.
        \item Filter candidates: must NOT be an end voxel, must be a simple point, must NOT have $\text{bc\_tag} > 0$.
        \item Sequential deletion with re-checking: after each deletion, re-verify remaining candidates (the ``double verification'' is essential because earlier deletions change neighbourhoods).
    \end{itemize}
    \item Iterate until no more voxels are deleted.
\end{enumerate}

Convergence is guaranteed by monotonic voxel removal. Typical convergence time is 8--12 iterations for domains under $100^3$.

\subsubsection{Topology Preservation}

A voxel is \emph{simple} if its removal does not change the topology of either the foreground or background:
\begin{itemize}
    \item Foreground test: exactly 1 connected component in 26-neighbourhood after removal.
    \item Background test: exactly 1 connected component in 6-neighbourhood (face-adjacent) of background voxels.
\end{itemize}

End voxels (those with exactly 1 foreground 26-neighbour) are never deleted---this prevents skeleton retraction. The Euler characteristic is invariant, so thinning is homotopic and maintains topology.

\subsubsection{Implementation: Numba JIT Acceleration}

All inner loops (\texttt{is\_simple\_point}, \texttt{get\_components\_26}, \texttt{get\_components\_6\_bg}, \texttt{find\_candidates}, \texttt{sequential\_delete}) are compiled with \texttt{@njit} for performance. The thinning stage is the second most expensive after Top3D.

\fig{2.1: Before/after thinning---$150 \times 40 \times 4$ cantilever voxel field reducing to skeleton (from history snapshots ``1\_Initial\_Voxels'' and ``2\_Skeleton\_Voxels'')}

\subsection{Graph Construction and Post-Processing}
\label{subsec:graph}

\subsubsection{Voxel Classification and Feature Detection}

Each skeleton voxel is classified by 26-neighbour count:
\begin{itemize}
    \item End: 1 neighbour
    \item Regular: 2 neighbours
    \item Joint: $> 2$ neighbours
\end{itemize}

End and Joint voxels are ``feature voxels''---they become graph nodes.

\subsubsection{Boundary Condition Cluster Consolidation}

BC-tagged voxels often form clusters. Connected component analysis on tagged clusters replaces each cluster with its centroid. Centroids are reconnected to the skeleton via 3D Bresenham line drawing. This prevents spurious graph fragmentation at load/support introduction points.

\subsubsection{Edge Tracing and Graph Assembly}

From each feature voxel, march through regular voxels until another feature voxel is reached. Each march produces one graph edge with intermediate polyline points. Laplacian smoothing is applied to intermediates to reduce staircase artefacts from the voxel grid. The result is a graph $G = (V, E)$ where $V$ are nodes (with $xyz$ positions, tags, degrees) and $E$ are edges (with polyline intermediates).

\subsubsection{Post-Processing Pipeline}

\begin{enumerate}
    \item \textbf{Edge collapse:} Merge nodes connected by edges shorter than a threshold; tagged nodes are protected from deletion.
    \item \textbf{Branch pruning:} Iteratively remove degree-1 (dangling) branches shorter than a minimum length; never prune branches that terminate at tagged nodes.
    \item \textbf{RDP geometric simplification:} Ramer--Douglas--Peucker algorithm applied to intermediate polyline points.
    \item \textbf{Radius estimation:} Two modes: (a) EDT-based: $\text{radius} = 2 \times \text{EDT}_{\text{value}} \times \text{pitch}$; (b) Uniform volume-matching: $r = \sqrt{V_{\text{target}} / (\pi L_{\text{total}})}$ (preferred because it preserves total volume).
    \item \textbf{Extrema anchoring:} Ensure graph nodes exist at bounding box limits.
\end{enumerate}

\alg{2: Pseudocode for graph extraction with BC consolidation}

\fig{2.2: Graph extraction stages side-by-side (4-panel: Raw Graph $\rightarrow$ After Collapse $\rightarrow$ After Prune $\rightarrow$ After RDP Simplification)}

\subsection{Frame Optimisation}
\label{subsec:frame_opt}

\subsubsection{3D Euler-Bernoulli Beam Finite Element Analysis}

Each beam element has 2 nodes with 6 DOF per node (3 translational + 3 rotational), totalling 12 DOF. For circular cross-section:
\begin{align}
    A &= \pi r^2 \\
    I_y = I_z &= \frac{\pi r^4}{4} \\
    J &= \frac{\pi r^4}{2}
\end{align}

The $12 \times 12$ local stiffness matrix incorporates axial, bending (two planes), and torsional stiffness. A rotation matrix transforms from local beam coordinates to global coordinates (special handling for vertical alignment). Global assembly uses sparse COO format, converted to CSC, solved via \texttt{scipy.sparse.linalg.spsolve}. Compliance is $C = \mathbf{u}^T \mathbf{K} \mathbf{u} = \mathbf{F}^T \mathbf{u}$.

\subsubsection{Size Optimisation (Optimality Criteria Method)}

\textbf{Design variables:} Beam radii $r_i$ (one per edge).

\textbf{Objective:} Minimise compliance $C = \mathbf{u}^T \mathbf{K} \mathbf{u}$.

\textbf{Constraint:} $\sum_i \pi r_i^2 L_i = V_{\text{target}}$.

\textbf{Sensitivities:}
\begin{equation}
    \frac{\partial C}{\partial r_i} = -\mathbf{u}_e^T \frac{\partial \mathbf{K}_e}{\partial r_i} \mathbf{u}_e,
\end{equation}
where $\frac{\partial \mathbf{K}}{\partial r}$ depends on $\frac{\partial A}{\partial r} = 2\pi r$, $\frac{\partial I}{\partial r} = \pi r^3$, $\frac{\partial J}{\partial r} = 2\pi r^3$.

The OC update rule with Lagrange multiplier bisection and move limits typically converges in 15--25 iterations. No minimum gauge constraint is enforced (future work).

\subsubsection{Layout Optimisation (L-BFGS-B)}

\textbf{Design variables:} $xyz$ coordinates of free (non-tagged) nodes.

\textbf{Fixed nodes:} Tagged nodes have locked coordinates.

\textbf{Bounds:} Node movement is limited to a move limit plus the design domain envelope offset by maximum radius.

\textbf{Method:} L-BFGS-B via \texttt{scipy.optimize.minimize} with finite differencing. Post-optimisation node snapping via Union-Find merges nodes within snapping distance. Layout optimisation changes topology by moving/merging nodes.

\subsubsection{Iterative Size-then-Layout Convergence}

\textbf{Key empirical finding:} Order matters. Size $\rightarrow$ Layout produces monotonic compliance reduction; Layout $\rightarrow$ Size causes compliance spikes (discovered during development, commit \texttt{0afff4c}).

\textbf{Convergence pattern:}
\begin{itemize}
    \item Iteration 1: Size optimisation achieves $\approx 50\%$ compliance reduction; Layout optimisation achieves additional $\approx 20\%$.
    \item Iteration 2: Size optimisation fine-tunes ($\approx 3$--5\% improvement); Layout optimisation provides marginal further improvement.
    \item Iterations 3+: Diminishing returns ($< 1\%$ improvement).
\end{itemize}

\textbf{Conclusion:} 2 iterations is sufficient. This is a practical finding with engineering significance---saves computation without sacrificing quality. No prior work (Yin 2020, Polak 2023) reports this convergence behaviour.

\tbl{2.2: Iterative convergence data---compliance at each stage across 1, 2, and 5 iteration loops, showing diminishing returns}

\fig{2.3: Compliance convergence plot ($x$-axis: pipeline stage, $y$-axis: compliance, multiple lines for different iteration counts)}

\subsection{Hybrid Beam-Plate Approach}
\label{subsec:hybrid}

\subsubsection{Motivation}

Topology optimisation results for problems with distributed loads (roofs, panels, shells) produce plate-like regions alongside beam-like members. Representing a plate as a network of beams is structurally inaccurate and geometrically poor. Automated classification of regions into beams vs.\ plates \emph{before} thinning allows different thinning modes to be applied.

\subsubsection{Zone Classification (Multi-Signal Scoring)}

Operates on connected components of the binarised voxel field. Four weighted signals computed per region:

\begin{enumerate}
    \item \textbf{PCA Shape (weight 0.3):} Compute eigenvalues $\lambda_1 \geq \lambda_2 \geq \lambda_3$ from PCA of voxel coordinates.
    \begin{align}
        \text{Planarity} &= 1 - \frac{\lambda_3}{\lambda_2} \quad \text{(high for plates)}, \\
        \text{Linearity} &= 1 - \frac{\lambda_2}{\lambda_1} \quad \text{(high for beams)}.
    \end{align}

    \item \textbf{Aspect Ratio (weight 0.3):} $\frac{\min(\lambda_1, \lambda_2)}{\max(\lambda_1, \lambda_2)}$---plates score $\approx 1$, beams score $\ll 1$.

    \item \textbf{BC Tag Density (weight 0.2):} Fraction of voxels with load tags biases toward plate; fraction with support tags biases toward beam.

    \item \textbf{EDT Uniformity (weight 0.2):} Coefficient of variation of EDT values within the region---plates uniform (low CoV), beams taper (high CoV).
\end{enumerate}

Final score: $\text{plate\_score} = \text{weighted sum}$; $\text{beam\_score} = 1 - \text{plate\_score}$. Threshold: $\text{plate\_score} > 0.5 \rightarrow \text{plate zone} (\text{zone} = 1)$, else $\text{beam zone} (\text{zone} = 2)$.

For large regions ($> 2000$ voxels): block-level analysis using $12^3$ sub-blocks with overall score as tiebreaker.

\subsubsection{Separate Thinning Modes}

\begin{itemize}
    \item Beam zones (zone = 2): thinned with mode = 0 (curve-preserving)---standard Yin thinning producing 1D medial curves.
    \item Plate zones (zone = 1): thinned with mode = 1 (surface-preserving)---modified thinning preserving 2D medial surfaces.
\end{itemize}

Interface boundary masking: voxels at zone boundaries are protected during thinning to prevent skeleton retraction at interfaces. \textbf{Key design principle:} Classify \emph{before} thinning---thinning destroys the geometric information (thickness, shape) needed for classification.

\subsubsection{Plate Extraction}

Dual output per plate region:
\begin{itemize}
    \item Solid mesh: Marching cubes on plate zone voxels for visualisation/CAD.
    \item Mid-surface with thickness: Each vertex gets $\text{thickness} = 2 \times \text{EDT}_{\text{value}} \times \text{pitch}$.
\end{itemize}

Mid-surface generation for FreeCAD export is currently limited (requires Open3D; noted as future work).

\subsubsection{Joint Creation}

Beam endpoints from skeleton extraction often don't coincide with plate boundary vertices. Solution: KD-tree spatial index on plate boundary vertices; snap each beam endpoint within a tolerance to the nearest plate boundary vertex. This creates structural connectivity between beam and plate elements in the final model.

\tbl{2.3: Zone classification results---plate\_score, beam\_score, classification, voxel counts for Clear\_Beam\_Plate, Thick\_Beam, Cantilever, Hybrid\_Roof}

\fig{2.4: Hybrid\_Roof zone classification visualisation (cyan = plate zones, red = beam zones)}

\subsection{CAD Model Generation}
\label{subsec:cad}

\subsubsection{CSG Tree Construction}

Beam members: cylinders (straight beams) or Bézier-swept tubes (curved beams) connecting node pairs. Joints: spheres at graph nodes with radius matching the maximum connected beam radius. Boolean union of all primitives into a single solid.

\subsubsection{Cubic Bézier Curve Fitting}

For edges with intermediate polyline points (curved skeleton paths):
\begin{itemize}
    \item Least-squares fit of interior control points $P_1$, $P_2$ given endpoints $P_0$, $P_3$ and intermediate samples.
    \item Sanitisation: chord monotonicity constraint ($0 \leq t_1 \leq t_2 \leq \text{chord\_length}$) and perpendicular bulge limit ($\text{max\_bulge\_ratio} \times \text{chord\_length}$) to prevent self-intersecting curves.
    \item Fallback: straight chord initialisation for edges with insufficient intermediates or failed fitting.
\end{itemize}

Curved beams improve geometric fidelity without significant compliance change.

\subsubsection{FreeCAD Macro Export}

Automated Python macro generation for FreeCAD:
\begin{itemize}
    \item \texttt{Part.BezierCurve} for curved beam centreline paths.
    \item \texttt{makePipeShell} with Frenet mode for swept tube geometry along Bézier paths.
    \item Ball-and-stick fallback (cylinder + sphere) for straight beams.
    \item Plate reconstruction: B-spline surfaces or batched face creation from mesh triangles.
\end{itemize}

Crash resistance: every geometry operation wrapped in \texttt{try/except} with graceful degradation (skip failed elements, log with traceback). Batch processing: \texttt{FreeCADGui.updateGui()} called every 10--20 iterations to prevent GUI freezing. Pipeline history visualisation: optional rendering of all intermediate stages (voxel snapshots, graph stages) within FreeCAD.

\fig{2.5: FreeCAD screenshot---reconstructed cantilever with curved Bézier beam members}

\section{Case Studies and Results}
\label{sec:results}

\subsection{Test Problems and Setup}
\label{subsec:setup}

Three primary test cases chosen to demonstrate different pipeline capabilities:

\begin{enumerate}
    \item \textbf{$150 \times 40 \times 4$ Cantilever} (beam-only): primary benchmark matching Yin's test case---tip load, fixed left face, \texttt{volfrac} = 0.3.
    \item \textbf{$60 \times 20 \times 2$ Short Cantilever} (beam-only): fast validation case with simpler topology.
    \item \textbf{$40 \times 40 \times 20$ Roof Structure} (hybrid beam-plate): distributed top load, corner supports---produces clear plate region with beam supports.
\end{enumerate}

Hardware and software: Python 3.12, NumPy, SciPy, Numba, Open3D, FreeCAD 0.21+. All results are reproducible via \texttt{run\_pipeline.py} with documented command-line arguments.

\tbl{3.1: Test problem parameters---domain dimensions, volfrac, load type/location, support type/location, expected structural type}

\subsection{Beam-Only Reconstruction Results}
\label{subsec:beam_results}

\subsubsection{$150 \times 40 \times 4$ Cantilever---Full Pipeline Walkthrough}

\begin{itemize}
    \item \textbf{Stage 0:} Top3D converges in $\approx 50$ iterations at \texttt{volfrac} = 0.3, producing a density field with clear truss-like topology.
    \item \textbf{Stage 1:} Reconstruction produces $\approx 15$ nodes and $\approx 18$ edges; skeleton preserves all BC points; graph post-processing removes short stubs and simplifies geometry.
    \item \textbf{Stage 3 (Size opt, iteration 1):} Baseline compliance $\approx 141,700$ reduced to $\approx 58,000$---structural members resized for load paths.
    \item \textbf{Stage 2 (Layout opt, iteration 1):} Compliance further reduced to $\approx 52,000$---nodes repositioned for structural efficiency.
    \item \textbf{Iteration 2:} Size $\rightarrow \approx 50,500$, Layout $\rightarrow \approx 49,000$---diminishing returns clearly visible.
    \item \textbf{Volume constraint:} $< 0.5\%$ error at every stage.
\end{itemize}

\subsubsection{Curved vs.\ Straight Beam Comparison}

Cubic Bézier fitting captures skeleton curvature faithfully. Compliance difference between curved and straight models is negligible ($< 2\%$). Visual quality improvement is significant---smoother, more organic geometry in FreeCAD.

\subsubsection{$60 \times 20 \times 2$ Short Cantilever}

Simpler topology (fewer members), faster convergence, validates pipeline on smaller domain.

\tbl{3.2: Full iterative optimisation comparison table---baseline, iter1-size, iter1-layout, iter2-size, iter2-layout compliance values, \% reduction, volume error}

\fig{3.1: $150 \times 40 \times 4$ cantilever pipeline stages (5-panel composite): (a) topology optimisation density field, (b) thinned skeleton, (c) extracted graph, (d) size+layout optimised frame, (e) FreeCAD CSG model}

\fig{3.2: Compliance convergence plot---compliance vs.\ pipeline stage for \texttt{opt\_loops} = 1, 2, and 5, showing 2-iteration convergence}

\subsection{Hybrid Beam-Plate Reconstruction Results}
\label{subsec:hybrid_results}

\subsubsection{$40 \times 40 \times 20$ Roof Structure}

Topology optimisation produces a horizontal plate region at top with vertical column beams at corners---the canonical beam-plate mixed structure. Zone classification: \texttt{plate\_score} = 0.65, \texttt{beam\_score} = 0.35 $\rightarrow$ correctly identified as mixed. Separate thinning: beam zones produce 1D medial curves, plate zone preserves 2D surface. Plate extraction: solid mesh generated successfully; mid-surface with per-vertex thickness computed. Joint creation: beam endpoints snapped to plate boundary within tolerance.

\subsubsection{Classification Accuracy}

\begin{itemize}
    \item Clear\_Beam\_Plate: 0.93 plate / 0.07 beam $\rightarrow$ mixed (1970 plate, 938 beam voxels)---correct.
    \item Thick\_Beam: 0.41 plate / 0.59 beam $\rightarrow$ beam (0 plate, 2932 beam voxels)---correct.
    \item Cantilever: 0.32 plate / 0.68 beam $\rightarrow$ beam (0 plate, 572 beam voxels)---correct.
    \item Hybrid\_Roof: 0.65 plate / 0.35 beam $\rightarrow$ mixed (8892 plate, 3892 beam voxels)---correct.
\end{itemize}

Multi-signal scoring correctly handles all cases including the ambiguous Thick\_Beam (which a PCA-only approach misclassified).

\fig{3.3: Roof structure zone classification visualisation (cyan = plate, red = beam, 3D voxel rendering)}

\tbl{3.3: Zone classification accuracy---plate\_score, beam\_score, classification, plate voxels, beam voxels for all 4 test cases}

\subsection{Comparison with Existing Approaches}
\label{subsec:comparison}

Systematic feature comparison across this work and all 4 reference papers:

\begin{itemize}
    \item \textbf{vs.\ Yin (2020):} Faithful reproduction of core pipeline + extensions (hybrid classification, iterative convergence, curved beams, Python, FreeCAD). Yin's MATLAB code is not publicly available; this work provides an open-source alternative.

    \item \textbf{vs.\ Amroune (2021):} Lofting approach produces smoother surfaces for individual branches, but cannot handle mixed beam-plate geometries and lacks structural re-optimisation.

    \item \textbf{vs.\ Feng (2025):} RMF approach provides editable STEP files with smooth blending surfaces, but requires complex RMF computation and performs no structural optimisation after reconstruction.

    \item \textbf{vs.\ Polak (2023):} Evolutionary refinement is conceptually similar to iterative optimisation but is heuristic-based with no formal convergence analysis; this work provides formal size+layout optimisation with demonstrated convergence.
\end{itemize}

Unique contributions: (1) only approach with demonstrated iterative convergence, (2) only approach with hybrid beam-plate handling, (3) only fully open-source Python implementation.

\tbl{3.4: Extended comparison table---columns: This Work, Yin 2020, Amroune 2021, Feng 2025, Polak 2023; rows: Language, Topology Optimisation method, Skeletonisation type, Element types, Zone classification, Structural re-optimisation, Iterative convergence, CAD format, Open source}

\section{Discussion}
\label{sec:discussion}

\subsection{Accuracy and Structural Performance}
\label{subsec:accuracy}

Total compliance reduction of 50--60\% across iterative optimisation demonstrates significant structural improvement over the raw skeleton-derived frame. The volume constraint is consistently satisfied within $< 0.5\%$---the OC bisection method is robust. There is a trade-off between structural performance and geometric fidelity: as compliance decreases, geometric likeness to the original topology optimisation result also decreases because nodes move away from medial axis positions for structural efficiency.

The 2-iteration convergence finding has practical engineering implications: designers can confidently run 2 iterations without diminishing returns, saving computation for large models. The compliance of the frame model relative to the original SIMP continuum compliance should be considered: the skeletal approximation introduces a compliance increase due to reduced geometric complexity, but the frame model is parametric and manufacturable.

\subsection{Robustness and Failure Modes}
\label{subsec:robustness}

\subsubsection{Thinning Stability}

Convergence is guaranteed by monotonic voxel removal. Numba JIT compilation handles domains up to $100^3$ efficiently. The double-verification loop prevents topological errors.

\subsubsection{Graph Extraction Edge Cases}

Isolated single voxels (no neighbours), branching cycles (resolved by consolidation), and degenerate zero-length edges (removed by collapse) are handled.

\subsubsection{Optimisation Stability}

Near-singular stiffness matrices when beams become very thin (small radius) are mitigated by minimum radius constraints. Layout optimisation can produce \texttt{NaN} displacements---caught and handled.

\subsubsection{FreeCAD Macro Robustness}

\texttt{shell.removeSplitter()} and \texttt{shell.makeOffsetShape()} are identified as unstable Part API calls and avoided in the current implementation. Every geometry creation is wrapped in \texttt{try/except} with graceful skip-and-log. Batch \texttt{updateGui()} prevents GUI freeze on large models.

\subsection{Limitations and Trade-offs}
\label{subsec:limitations}

\begin{itemize}
    \item \textbf{No shell FEA for plates:} Plate regions are classified and extracted but not included in the frame FEA model---no shell elements are implemented. The frame optimisation operates only on beam members. This limits the structural validation of hybrid models.

    \item \textbf{Mid-surface generation fragility:} Currently skipped in the FreeCAD macro because Open3D-based Poisson reconstruction is unreliable for thin plate geometries.

    \item \textbf{Fixed block size:} Zone classifier uses $12^3$ sub-blocks for large regions---not adaptive and may need tuning for very large or very small domains.

    \item \textbf{Single load case:} Current implementation handles one load case; multi-load-case optimisation requires extension of the FEA solver and OC formulation.

    \item \textbf{No manufacturing constraints:} Pipeline produces geometrically valid CAD but does not enforce overhang angles, minimum feature sizes, or build direction constraints for additive manufacturing.

    \item \textbf{Computational cost:} Top3D dominates runtime for large domains; thinning is second; optimisation is relatively fast due to low DOF count of the frame model.
\end{itemize}

\section{Conclusions and Future Work}
\label{sec:conclusions}

\subsection{Summary of Contributions}
\label{subsec:summary}

This dissertation presents four contributions with specific evidence from case studies:

\begin{enumerate}
    \item \textbf{Python pipeline} validated against Yin's approach on identical test cases (150$\times$40$\times$4 cantilever).

    \item \textbf{Multi-signal zone classification} correctly classifies all four test structures, including ambiguous cases that single-metric approaches fail on.

    \item \textbf{Iterative size-then-layout optimisation} converges in 2 iterations with 50--60\% total compliance reduction; ordering (Size THEN Layout) is critical.

    \item \textbf{End-to-end automation} produces parametric FreeCAD models with curved Bézier beams from raw SIMP density fields.
\end{enumerate}

This is a complete, working, open-source system---not a theoretical proposal or partial prototype.

\subsection{Directions for Future Research}
\label{subsec:future}

\begin{itemize}
    \item \textbf{Shell finite elements for plates:} Integrate shell element formulation to include plate regions in structural optimisation, enabling true hybrid beam-shell optimisation.

    \item \textbf{Multi-load-case optimisation:} Extend FEA and OC formulation to handle multiple load cases simultaneously (weighted compliance sum).

    \item \textbf{Machine learning for zone classification:} Train a classifier on labelled zone data to replace or augment the hand-tuned signal weights.

    \item \textbf{Additive manufacturing constraints:} Incorporate overhang angle limits, minimum feature size, and build direction into the optimisation formulation.

    \item \textbf{Commercial CAD integration:} Extend export beyond FreeCAD to STEP/IGES for SolidWorks, Fusion 360, etc.

    \item \textbf{Large-domain scalability:} Profile and optimise the zone classifier and thinning stages for domains exceeding $100^3$ elements.
\end{itemize}

\end{document}
